\editionPage{defi-ex-1}{Défi D01 · Deux laboratoires}{Chapitre 1 · Approfondissement}{Chercher, organiser, justifier · 25 à 40 minutes}
\label{defi-ex-1-1}\label{defis-debut}
Les laboratoires A et B mesurent la même longueur. A : $\bar L=\qty{50,12}{mm}$, $s=\qty{0,08}{mm}$, $N=16$. B : $\bar L=\qty{50,10}{mm}$, $s=\qty{0,12}{mm}$, $N=36$. La référence est \qty{50,00}{mm}.
\begin{questions}
\item Calculer les deux incertitudes-types sur la moyenne.
\item Comparer les dispersions individuelles puis les incertitudes sur la moyenne.
\item Que suggère le décalage commun par rapport à la référence ?
\end{questions}
\smallskip\otRenvoi{Indications}{indice-1}\hfill\otRenvoi{Corrigé}{defi-cor-1-1}

\editionPage{defi-ex-2}{Défi D02 · Contrôler une étiquette}{Chapitre 2 · Approfondissement}{Chercher, organiser, justifier · 25 à 40 minutes}
\label{defi-ex-2-1}
Une bouteille annoncée à 15,0 \% de sucre en masse contient 250 mL de solution. La bouteille vide pèse 32,0 g et pleine 302,0 g. Une analyse indépendante trouve 40,5 g de sucre.
\begin{questions}
\item Déterminer la masse volumique de la solution.
\item Vérifier le pourcentage annoncé.
\item Expliquer pourquoi « 15 g pour 100 mL » serait une autre annonce.
\end{questions}
\smallskip\otRenvoi{Indications}{indice-2}\hfill\otRenvoi{Corrigé}{defi-cor-2-1}

\editionPage{defi-ex-3}{Défi D03 · Préparer avec le stock disponible}{Chapitre 3 · Approfondissement}{Chercher, organiser, justifier · 25 à 40 minutes}
\label{defi-ex-3-1}
On veut 250,0 mL à \qty{6,00}{g.L^{-1}} avec une mère à \qty{30,0}{g.L^{-1}}. On dispose de pipettes de 10,0 et 25,0 mL et d’une fiole de 250,0 mL.
\begin{questions}
\item Calculer le volume mère nécessaire.
\item Proposer un protocole utilisant le matériel disponible.
\item Une prise de 25,0 mL complétée à 250,0 mL conviendrait-elle ?
\end{questions}
\smallskip\otRenvoi{Indications}{indice-3}\hfill\otRenvoi{Corrigé}{defi-cor-3-1}

\editionPage{defi-ex-4}{Défi D04 · Reconstituer une entité}{Chapitre 4 · Approfondissement}{Chercher, organiser, justifier · 25 à 40 minutes}
\label{defi-ex-4-1}
Un noyau possède 16 nucléons dont 8 neutrons. L’entité qui le contient porte la charge $-2e$.
\begin{questions}
\item Identifier l’élément et écrire le noyau.
\item Déterminer le nombre d’électrons.
\item Associer cet ion à $\ce{Na+}$ pour écrire un composé neutre.
\end{questions}
\smallskip\otRenvoi{Indications}{indice-4}\hfill\otRenvoi{Corrigé}{defi-cor-4-1}

\editionPage{defi-ex-5}{Défi D05 · Retrouver un élément}{Chapitre 5 · Approfondissement}{Chercher, organiser, justifier · 25 à 40 minutes}
\label{defi-ex-5-1}
Un atome appartient à la période 3 et possède sept électrons de valence. On se limite aux 18 premiers éléments.
\begin{questions}
\item Écrire sa configuration et identifier l’élément.
\item Prévoir la charge de son ion usuel.
\item Expliquer pourquoi l’ion et le gaz noble voisin restent deux espèces différentes.
\end{questions}
\smallskip\otRenvoi{Indications}{indice-5}\hfill\otRenvoi{Corrigé}{defi-cor-5-1}

\editionPage{defi-ex-6}{Défi D06 · Une solution vue de deux échelles}{Chapitre 6 · Approfondissement}{Chercher, organiser, justifier · 25 à 40 minutes}
\label{defi-ex-6-1}
Une solution de volume 250 mL a une concentration de glucose de \qty{12,0}{g.L^{-1}}. On donne $m_C=1,99\times10^{-26}$ kg, $m_H=1,67\times10^{-27}$ kg, $m_O=2,66\times10^{-26}$ kg et $N_A=6,02\times10^{23}\unit{mol^{-1}}$.
\begin{questions}
\item Calculer la masse d’une molécule \ce{C6H12O6}.
\item Déterminer le nombre de molécules de glucose dans la solution.
\item En déduire sa quantité de matière.
\end{questions}
\smallskip\otRenvoi{Indications}{indice-6}\hfill\otRenvoi{Corrigé}{defi-cor-6-1}

\editionPage{defi-ex-7}{Défi D07 · Une réserve de froid}{Chapitre 7 · Approfondissement}{Chercher, organiser, justifier · 25 à 40 minutes}
\label{defi-ex-7-1}
Un accumulateur contient 0,300 kg de glace à 0 °C. Le milieu lui transmet en moyenne 25 J chaque seconde. $L_f=\qty{3,34e5}{J.kg^{-1}}$.
\begin{questions}
\item Déterminer l’énergie absorbable avant fusion complète.
\item Estimer la durée de maintien du mélange glace-eau à 0 °C.
\item Expliquer ce que le modèle ne prévoit pas après cette durée.
\end{questions}
\smallskip\otRenvoi{Indications}{indice-7}\hfill\otRenvoi{Corrigé}{defi-cor-7-1}

\editionPage{defi-ex-8}{Défi D08 · Choisir le bon mélange}{Chapitre 8 · Approfondissement}{Chercher, organiser, justifier · 25 à 40 minutes}
\label{defi-ex-8-1}
On veut produire exactement 0,300 mol d’eau par $\ce{2 H2 + O2 -> 2 H2O}$, sans réactif restant.
\begin{questions}
\item Déterminer les quantités initiales nécessaires.
\item On introduit en réalité 0,300 mol de chaque gaz : calculer l’état final.
\item Comparer la quantité d’eau obtenue dans les deux cas.
\end{questions}
\smallskip\otRenvoi{Indications}{indice-8}\hfill\otRenvoi{Corrigé}{defi-cor-8-1}

\editionPage{defi-ex-9}{Défi D09 · Le produit est-il prêt ?}{Chapitre 9 · Approfondissement}{Chercher, organiser, justifier · 25 à 40 minutes}
\label{defi-ex-9-1}
La CCM du produit synthétisé S présente une tache correspondant au produit naturel N et une autre correspondant au réactif R.
\schema{ccm-synthese}
\begin{questions}
\item Peut-on conclure que la synthèse a formé l’espèce recherchée ?
\item Peut-on livrer S comme espèce pure sur cette seule analyse ?
\item Proposer une stratégie de contrôle après purification.
\end{questions}
\smallskip\otRenvoi{Indications}{indice-9}\hfill\otRenvoi{Corrigé}{defi-cor-9-1}

\editionPage{defi-ex-10}{Défi D10 · Auditer une réaction}{Chapitre 10 · Approfondissement}{Chercher, organiser, justifier · 25 à 40 minutes}
\label{defi-ex-10-1}
On propose $\ce{^{235}_{92}U + ^{1}_{0}n -> ^{140}_{54}Xe + ^{94}_{38}Sr + x ^{1}_{0}n}$. Chaque événement libère ici $3,0\times10^{-11}$ J, valeur fournie pour l’exercice.
\begin{questions}
\item Trouver $x$ et vérifier la charge.
\item Nommer la famille de transformation.
\item Combien d’événements faut-il pour libérer \qty{9,0e7}{J} ?
\end{questions}
\smallskip\otRenvoi{Indications}{indice-10}\hfill\otRenvoi{Corrigé}{defi-cor-10-1}

\editionPage{defi-ex-11}{Défi D11 · Valider un pointage}{Chapitre 11 · Approfondissement}{Chercher, organiser, justifier · 25 à 40 minutes}
\label{defi-ex-11-1}
Une vidéo fournit les positions de la chronophotographie. Le pas vaut 0,20 s et les abscisses 0 ; 0,20 ; 0,80 ; 1,80 ; 3,20 m.
\schema{chronophoto}
\begin{questions}
\item Calculer les vitesses approchées sur les quatre intervalles.
\item Représenter $\vect v_1$ et $\vect v_3$ à l’échelle 1 cm pour 2 m/s.
\item Peut-on conclure que la vitesse augmente de façon exactement régulière entre les images ?
\end{questions}
\smallskip\otRenvoi{Indications}{indice-11}\hfill\otRenvoi{Corrigé}{defi-cor-11-1}

\editionPage{defi-ex-12}{Défi D12 · Une pesanteur plus faible}{Chapitre 12 · Approfondissement}{Chercher, organiser, justifier · 25 à 40 minutes}
\label{defi-ex-12-1}
Un astronaute de 80 kg se trouve à une distance du centre terrestre égale à deux rayons terrestres. Au sol, on adopte $g_0=9,8$ N/kg.
\begin{questions}
\item Comparer l’attraction à celle au sol.
\item Calculer son poids gravitationnel à cette distance.
\item Un poids nul est-il nécessaire pour être en orbite ?
\end{questions}
\smallskip\otRenvoi{Indications}{indice-12}\hfill\otRenvoi{Corrigé}{defi-cor-12-1}

\editionPage{defi-ex-13}{Défi D13 · Trois phases d’un lancer}{Chapitre 13 · Approfondissement}{Chercher, organiser, justifier · 25 à 40 minutes}
\label{defi-ex-13-1}
Une balle est lancée vers le haut puis retombe ; on utilise le modèle de chute libre. La figure représente trois instants : montée, sommet, descente.
\schema{chute}
\begin{questions}
\item Compléter les directions et sens de la vitesse et du poids à chaque phase.
\item Expliquer pourquoi la balle n’est pas en équilibre au sommet.
\item Comment tester la pertinence du modèle avec une vidéo ?
\end{questions}
\smallskip\otRenvoi{Indications}{indice-13}\hfill\otRenvoi{Corrigé}{defi-cor-13-1}

\editionPage{defi-ex-14}{Défi D14 · Une galerie à mesurer}{Chapitre 14 · Approfondissement}{Chercher, organiser, justifier · 25 à 40 minutes}
\label{defi-ex-14-1}
Un signal bref est émis à l’entrée d’une galerie rectiligne. Un écho arrive 0,420 s plus tard. On adopte 340 m/s. Une seconde mesure donne 0,424 s.
\schema{echo}
\begin{questions}
\item Estimer la longueur de la galerie avec chaque mesure.
\item Comparer les résultats et citer une cause de différence.
\item Comment vérifier que l’écho étudié provient bien du fond ?
\end{questions}
\smallskip\otRenvoi{Indications}{indice-14}\hfill\otRenvoi{Corrigé}{defi-cor-14-1}

\editionPage{defi-ex-15}{Défi D15 · Une note mal identifiée}{Chapitre 15 · Approfondissement}{Chercher, organiser, justifier · 25 à 40 minutes}
\label{defi-ex-15-1}
Un enregistrement montre quatre périodes complètes en 9,09 ms. Un élève mesure seulement la durée entre un maximum et le minimum suivant, puis annonce une fréquence deux fois trop grande.
\begin{questions}
\item Calculer la fréquence à partir des quatre périodes.
\item Expliquer l’erreur de l’élève.
\item Un autre instrument joue avec la même période et une autre forme : que peut-on comparer ?
\end{questions}
\smallskip\otRenvoi{Indications}{indice-15}\hfill\otRenvoi{Corrigé}{defi-cor-15-1}

\editionPage{defi-ex-16}{Défi D16 · Identifier un éclairage}{Chapitre 16 · Approfondissement}{Chercher, organiser, justifier · 25 à 40 minutes}
\label{defi-ex-16-1}
Une lampe présente trois raies à 450, 520 et 650 nm sur un faible fond continu. Une autre lampe produit un spectre continu plus riche vers le bleu lorsqu’on augmente son chauffage.
\begin{questions}
\item Comparer les deux mécanismes d’émission proposés par les observations.
\item La première lampe est-elle monochromatique ?
\item Quelle information supplémentaire permettrait d’identifier l’espèce responsable des raies ?
\end{questions}
\smallskip\otRenvoi{Indications}{indice-16}\hfill\otRenvoi{Corrigé}{defi-cor-16-1}

\editionPage{defi-ex-17}{Défi D17 · Deux couleurs dans le verre}{Chapitre 17 · Approfondissement}{Chercher, organiser, justifier · 25 à 40 minutes}
\label{defi-ex-17-1}
Un faisceau arrive de l’air sous $i=40,0^\circ$. Le verre a pour indices $n_R=1,50$ pour le rouge et $n_V=1,53$ pour le violet.
\begin{questions}
\item Calculer les deux angles réfractés.
\item Quelle radiation se rapproche le plus de la normale à cette interface ?
\item Pourquoi une différence faible peut-elle devenir visible après le prisme ?
\end{questions}
\smallskip\otRenvoi{Indications}{indice-17}\hfill\otRenvoi{Corrigé}{defi-cor-17-1}

\editionPage{defi-ex-18}{Défi D18 · Agrandir une projection}{Chapitre 18 · Approfondissement}{Chercher, organiser, justifier · 25 à 40 minutes}
\label{defi-ex-18-1}
Un objet de 2,0 cm est placé à 15 cm d’une lentille de focale 10 cm. On souhaite le projeter sur un écran.
\schema{lentille-agrandie}
\begin{questions}
\item Construire l’image et déterminer la position de l’écran.
\item Calculer le grandissement par les tailles puis vérifier par Thalès.
\item Comparer à l’objet placé à 30 cm de la même lentille : quel choix agrandit l’image et quelle place faut-il prévoir derrière la lentille ?
\end{questions}
\smallskip\otRenvoi{Indications}{indice-18}\hfill\otRenvoi{Corrigé}{defi-cor-18-1}

\editionPage{defi-ex-19}{Défi D19 · Un circuit à diagnostiquer}{Chapitre 19 · Approfondissement}{Chercher, organiser, justifier · 25 à 40 minutes}
\label{defi-ex-19-1}
Une alimentation idéale de 6,0 V dessert deux branches en dérivation : $R_1=200\,\Omega$ et $R_2=300\,\Omega$. Un ampèremètre principal indique 50 mA.
\schema{parallele-defi}
\begin{questions}
\item Calculer les intensités attendues dans les deux branches.
\item Comparer l’intensité totale prévue à la mesure.
\item Un élève propose de placer l’ampèremètre directement sur le générateur pour confirmer : corriger sa proposition.
\end{questions}
\smallskip\otRenvoi{Indications}{indice-19}\hfill\otRenvoi{Corrigé}{defi-cor-19-1}

\editionPage{defi-ex-20}{Défi D20 · Du code à l’alerte}{Chapitre 20 · Approfondissement}{Chercher, organiser, justifier · 25 à 40 minutes}
\label{defi-ex-20-1}
Le montage utilise $U_0=5,00$ V et $R_0=10,0$ kΩ. Le convertisseur fictif donne 1000 unités pour 5,00 V. Le capteur suit $R_c(\mathrm{k}\Omega)=20,0-0,400\theta$ de 0 à 40 °C. Une acquisition renvoie 400.
\schema{pont-capteur}
\begin{questions}
\item Calculer U puis $R_c$.
\item Déterminer la température et l’état d’une alerte au-dessus de 30 °C.
\item Proposer deux garde-fous dans le programme.
\end{questions}
\smallskip\otRenvoi{Indications}{indice-20}\hfill\otRenvoi{Corrigé}{defi-cor-20-1}

